\begin{titlepage}
\centering
{\Large Notas didácticas de una entrevista a fondo\par}
\vspace{1.0cm}
{\huge\bfseries Gemini Robotics, transferencia entre encarnaciones y modelos del mundo para robots\par}
\vspace{0.5cm}
{\Large Zhang Xiaojun conversa con Tan Jie (谭杰), de Google DeepMind\par}
\vspace{0.35cm}
{\large Elaborado con base en el video público de Zhang Xiaojun Podcast, una transcripción generada con GPU, la estructura de capítulos y diagramas conceptuales propios\par}
\vspace{0.3cm}
{\large 2025-11-28\par}
\vspace{0.85cm}
\includegraphics[width=0.88\textwidth,height=0.42\textheight,keepaspectratio]{cover.jpg}\par
\vfill
\begin{tcolorbox}[width=0.92\textwidth, colback=black!2!white, colframe=black!55, sharp corners]
\textbf{Título del video}: 121. Entrevista a Tan Jie (谭捷), de DeepMind: robots, transferencia entre encarnaciones, modelos del mundo, Gemini Robotics 1.5 y Google\par
\textbf{Canal del video}: Zhang Xiaojun Podcast\par
\textbf{Duración del video}: 02:06:16\par
\textbf{Enlace del video}: \href{\videourl}{\nolinkurl{\videourl}}\par
\textbf{Descripción del material}: YouTube no tiene subtítulos disponibles; el texto se elaboró con base en la transcripción de faster-whisper large-v3 con CUDA, los capítulos del video, la portada y figuras didácticas propias. Las tomas fijas de la entrevista no se reutilizan en el texto.\par
\textbf{Página del proyecto}: \href{\repourl}{\nolinkurl{\repourl}}\par
\end{tcolorbox}
\end{titlepage}

\tableofcontents
\newpage

